\section{The “Prophecies” of Rene Guenon (2)}

It might seem odd that the concluding chapter on a book about Hindu doctrines is all about the ``West''. Obviously, the intent is not for the West to become Hindu, but rather for the West to use the traditional doctrines as revealed in the Vedanta in order to become what it truly is.

Guenon ends the chapter with a deeper review of the three alternatives he mentioned. ``Everything clearly depends on the mental state of the Western world at the moment when it reaches the furthest term of its present civilization.'' If the mental state remains as it currently is -- viz., based on non-traditional values which both Evola and Guenon have analysed in various books -- then the first alternative will prevail. This is because there would be nothing left in the Western mentality to accept Tradition, either from within, but also by way of assimilation.

The second alternative -- assimilation -- would still require an intellectual kernel, even if confined to a numerically small elect. This elect would have to be strong enough to serve as intermediaries to guide the West back to the sources of ``true intellectuality''. This does not require the conscious awareness of the masses to be effective. Guenon warns that this alternative ``would not be free from certain unpleasant features''. These unpleasant features are a matter for speculation, but probably not difficult to discern.

Now, if the elect were strong enough and had sufficient time to operate on its own, then the third possibility could be realised. It would have to act as a ``leaven'', slowly influencing the Western world. This does not require the consent of the mass of Westerners, nor even for the generality of so-called intellectuals. ``It is enough, therefore, as a start, for a few individuals to understand the need for such a change, but of course on condition that they understand it truly and thoroughly.''

Guenon points out: ``A return to a traditional civilization, both in principle and in respect of the whole body of institutions, is evidently the base condition for the transformation we have been speaking about.'' Clearly, a return to tradition is the most essential purpose of the elect, but how to bring this about. A concrete example would be helpful.

As an example, Guenon provides: ``We can only say that the Middle Ages afford us an example of a traditional development that was truly Western; ultimately it would be a case not purely and simply of copying or reconstructing what existed then, but of drawing inspiration from it in order to bring about an adaptation to suit the actual circumstances. \emph{If there exists a ``Western tradition'', that is where is must be looked for, and not in the fantasies of occultists and pseudo-esotericists.}'' My emphasis. Lest anyone thing Guenon is being idiosyncratic here, compare to Evola: ``the grandeur of Rome, having risen from the forces of the Nordic Aryans, created the last, great, universal period in the West, the feudal-imperial civilization of the Middle Ages.'' (from ``Pagan Imperialism'') And again, ``The power of a new Middle Ages is needed.'' Sir John Woodroffe also regards that era as traditional in comparison to the modern West. (``The World as Power'')

Unless there is a secret ``elect'' at work, why is it that no one looks to that model? Instead, we have the ``Nouvelle Droite'', Asatru and similar neo-pagan reconstructions, revivals of Germanic occultist organisation, and so on. The Middle Ages seem to be an embarrassment to everyone. To the contrary, to the extent that one cannot recognize the traditional elements at play there, to an even greater extent, it is unlikely that traditional elements can be recognized anywhere at all, never mind serve as a leaven for transformation.

\flrightit{Posted on 2008-01-14 by Cologero}
